\documentclass[../../main.tex]{subfiles}% 注意这里的文件路径不能用 ./main.tex ,否则用latexmk编译子文件会报错
\graphicspath{{\subfix{./image/}}{\subfix{../../image/}}} % 指定图片引用路径,后续可以直接使用图片文件名
% 如果图片存放在上级目录,则使用 ../../image/ 作为备用路径

% 例如:
% \begin{figure}[H]
% \centering
% \includegraphics[width=0.3\textwidth]{图.png}
% \caption{}
% \label{figure:图}
% \end{figure}
% 注意:上述\label{}一定要放在\caption{}之后,否则引用图片序号会只会显示??.

\begin{document}

\section{单调函数的可微性}

\subsection{Vitali覆盖定理}

\begin{definition}[Vitali覆盖]
设 $E\subset\mathbb{R}$, $\Gamma=\{I_\alpha\}$ 是一个区间族. 若对$\forall x\in E$,$\forall \varepsilon>0$, 存在 $I_\alpha\in\Gamma$, 使得 $x\in I_\alpha$, $|I_\alpha|<\varepsilon$, 则称 $\Gamma$ 是 $E$ 在 Vitali 意义下的一个覆盖, 简称为 $E$ 的 $\mathbf{Vitali}$\textbf{(维他利)覆盖}.
\end{definition}

\begin{example}
设 $E=[a,b]$, 记 $\{r_n\}$ 为 $[a,b]$ 中有理数的全体, 令
\begin{align*}
I_{n,m}=\left[r_n-\frac{1}{m}, r_n+\frac{1}{m}\right].
\end{align*}
则 $\mathcal{F}=\{I_{n,m}\}_{n,m=1}^\infty$ 是 $E$ 的Vitali覆盖.
\end{example}
\begin{proof}
任取 $x\in[a,b]$, 对 $\forall\varepsilon>0$, $\exists m_0\in\mathbb{N}$ 使得 $1/m_0<\varepsilon/2$. 对于 $1/m_0>0$, 由有理数的稠密性, 存在有理数 $r_{n_0}\in[a,b]$ 使得
\begin{align*}
|x-r_{n_0}|<\frac{1}{m_0}<\frac{\varepsilon}{2}.
\end{align*}
从而
\begin{align*}
x\in\left[r_{n_0}-\frac{1}{m_0}, r_{n_0}+\frac{1}{m_0}\right]\subset\left(r_{n_0}-\frac{\varepsilon}{2}, r_{n_0}+\frac{\varepsilon}{2}\right).
\end{align*}
记 $I_{n_0,m_0}=[r_{n_0}-1/m_0, r_{n_0}+1/m_0]$, 则 $x\in I_{n_0,m_0}$ 且 $|I_{n_0,m_0}|<\varepsilon$.
\end{proof}

\begin{theorem}[Vitali覆盖定理]\label{theorem:Vitali覆盖定理--实变函数}
设 $E\subset\mathbb{R}$, $m^*(E)<+\infty$, $\mathcal{F}$ 是 $E$ 的一个Vitali覆盖, 则对 $\forall\varepsilon>0$, 都存在有限个互不相交的区间 $I_1,\cdots,I_n\in\mathcal{F}$, 使得
\begin{align*}
m^*\!\left(E-\bigcup_{k=1}^n I_k\right)<\varepsilon.
\end{align*}
\end{theorem}
\begin{remark}
Vitali覆盖定理表明选出可列个互不相交的区间不一定能覆盖 $E$, 但覆盖不住的点集为一零测集.
\end{remark}
\begin{proof}
由于对任意的 $I_1,\cdots,I_n\in\mathcal{F}$, 它们端点的全体是零测集, 故不妨设 $\mathcal{F}$ 中的每个区间都是闭集.

由于 $m^*(E)<+\infty$, 由外测度定义, 存在开集 $G\supset E$ 使得 $m(G)<+\infty$. 不妨设 $\mathcal{F}$ 中的每个区间均包含在 $G$ 中(否则用 $\mathcal{F}'=\{I:I\in\mathcal{F}\ \text{且}\ I\subset G\}$ 代替 $\mathcal{F}$ 即可).
\begin{enumerate}[(1)]
\item 若存在有限个区间 $I_1,\cdots,I_n\in\mathcal{F}$, 使得 $E\subset\bigcup_{i=1}^n I_i$, 则 $m^*\!\left(E-\bigcup_{i=1}^n I_i\right)=0$, 定理结论成立.
\item 若对任意的 $I_1,\cdots,I_n\in\mathcal{F}$, 都有 $E\not\subset\bigcup_{i=1}^n I_i$. 在 $\mathcal{F}$ 中任取一个区间记为 $I_1$, 假设已选定互不相交的区间 $I_1,\cdots,I_k$, 由于 $E-\bigcup_{i=1}^k I_i\neq\varnothing $, 注意到 $\bigcup_{i=1}^k I_i$ 是闭集以及 $\mathcal{F}$ 是Vitali覆盖, 故至少存在一个 $I\in\mathcal{F}$, 使得 $I$ 与 $I_1,\cdots,I_k$ 都互不相交. 令
\begin{align*}
\lambda_k=\sup\{|I|:I\in\mathcal{F}, I\cap I_i=\varnothing , i=1,\cdots,k\}.
\end{align*}
注意到 $\mathcal{F}$ 中每个 $I$ 都含于 $G$ 中, 故 $\lambda_k\leqslant m(G)<+\infty$. 在 $\mathcal{F}$ 中选取区间 $I_{k+1}$, 使得
\begin{align}
|I_{k+1}|>\frac{\lambda_k}{2},\quad I_{k+1}\cap I_i=\varnothing ,\quad i=1,\cdots,k.\label{eq:::-hr82h39rhswufh8jiojshiojfw8290rui9ejwf90382jfwds1}
\end{align}
继续这一个过程, 可得到 $G$ 中一列互不相交的区间 $\{I_n\}$, 使得对每个 $k\in\mathbb{N}$ 都满足式 \eqref{eq:::-hr82h39rhswufh8jiojshiojfw8290rui9ejwf90382jfwds1}. 由于 $\bigcup_{k=1}^\infty I_k\subset G$, 故
\begin{align}
\sum_{k=1}^\infty|I_k|\leqslant m(G)<+\infty.\label{eq:::-hr82h39rhswufh8jiojshiojfw8290rui9ejwf90382jfwds2}
\end{align}
于是, 对 $\forall\varepsilon>0$, 存在 $n\in\mathbb{N}$ 使得 $\sum_{k=n+1}^\infty|I_k|<\varepsilon/5$. 令 $A=E-\bigcup_{k=1}^n I_k$, 下面证 $m^*(A)<\varepsilon$.

设 $x\in A$, 由于 $\bigcup_{k=1}^n I_k$ 是闭集且 $x\notin\bigcup_{k=1}^n I_k$. 又 $\mathcal{F}$ 是Vitali覆盖, 故存在 $I\in\mathcal{F}$, 使得 $I$ 包含 $x$ 且与 $I_1,\cdots,I_n$ 都不相交. 若进一步, $I$ 与 $\{I_k\}_{k>n}$ 中的每个区间都不相交, 则对任意的 $k>n$, 都有
\begin{align*}
|I|\leqslant\lambda_k<2|I_{k+1}|.
\end{align*}
再由式 \eqref{eq:::-hr82h39rhswufh8jiojshiojfw8290rui9ejwf90382jfwds2} 知, 当 $k\to\infty$ 时有 $|I_k|\to 0$, 从而 $|I|=0$, 于是 $I$ 是单点集, 矛盾. 因此 $I$ 必与 $\{I_k\}_{k>n}$ 中的某个区间相交. 令 $k_0=\min\{k:I\cap I_k\neq\varnothing \}$, 则 $k_0>n$, 且 $|I|\leqslant\lambda_{k_0-1}<2|I_{k_0}|$. 记 $I_{k_0}$ 的中心为 $x_{k_0}$, 半径为 $r_{k_0}$. 由于 $x\in I$ 且 $I\cap I_{k_0}\neq\varnothing $, 故
\begin{align*}
|x-x_{k_0}|\leqslant|I|+\frac{|I_{k_0}|}{2}<2|I_{k_0}|+\frac{|I_{k_0}|}{2}=\frac{5}{2}|I_{k_0}|.
\end{align*}
于是记 $J_{k_0}=[x_{k_0}-5r_{k_0},x_{k_0}+5r_{k_0}]$, 则 $x\in J_{k_0}$. 对每个 $I_k\in\{I_k\}_{k>n}$, 令 $J_k$ 为与 $I_k$ 具有相同中心且长度为 $I_k$ 的 $5$ 倍的区间, 则由前面的证明可知, $A\subset\bigcup_{k=n}^\infty J_k$. 因此
\begin{align*}
m^*(A)\leqslant\sum_{k=n+1}^\infty|J_k|=5\sum_{k=n+1}^\infty|I_k|<\varepsilon.
\end{align*}

\end{enumerate}
\end{proof}

\begin{definition}
设 $E\subset\mathbb{R}^n$, $\mathcal{F}$ 是 $\mathbb{R}^n$ 中的一族闭方体. 若对$\forall x\in E$,$\forall \varepsilon>0$, 存在 $Q\in\mathcal{F}$, 使得 $x\in Q$, 且有 $\mathrm{diam}Q<\varepsilon$, 则称 $\mathcal{F}$ 是 $E$ 的一个 $\mathbf{Vitali}$\textbf{覆盖}.
\end{definition}

\begin{proposition}
(Vitali) 设 $E\subset\mathbb{R}^n$ 是有界可测集, $\mathcal{F}$ 是 $E$ 的一个 Vitali 覆盖, 则存在互不相重(内部互不相交)的可数列 $\{Q_k\}$, 使得
\begin{align*}
m\!\left(E\setminus\bigcup_{k\geqslant 1} Q_k\right)=0.
\end{align*}
\end{proposition}
\begin{proof}

\end{proof}




\subsection{单调函数的可微性}

\begin{definition}[Dini导数]
设 $f(x)$ 是定义在 $\mathbb{R}$ 中点 $x_0$ 的一个邻域上的实值函数, 令
\begin{align*}
D^+f(x_0)=\mathop {\overline{\lim }} \limits_{h\rightarrow 0^+}\frac{f(x_0+h)-f(x_0)}{h},
\\
D_+f(x_0)=\mathop {\underline{\lim }} \limits_{h\rightarrow 0^+}\frac{f(x_0+h)-f(x_0)}{h},
\\
D^-f(x_0)=\mathop {\overline{\lim }} \limits_{h\rightarrow 0^-}\frac{f(x_0+h)-f(x_0)}{h},
\\
D_-f(x_0)=\mathop {\underline{\lim }} \limits_{h\rightarrow 0^-}\frac{f(x_0+h)-f(x_0)}{h},
\end{align*}
分别称它们为 $f(x)$ 在 $x_0$ 点的\textbf{右上导数}, \textbf{右下导数}, \textbf{左上导数}, \textbf{左下导数}. 总称为$\mathbf{Dini}$\textbf{(导)数}. 

此外,若此四个 Dini 数皆等于同一个有限值, 则 $f(x)$ 在 $x_0$ 点是\textbf{可微的}; 

若 $D^+f(x_0)=D_+f(x_0)$ 为有限值, 则 $f(x)$ 在 $x_0$ 点的\textbf{右导数存在}; 

若 $D^-f(x_0)=D_-f(x_0)$ 为有限值, 则 $f(x)$ 在 $x_0$ 点的\textbf{左导数存在}.
\end{definition}
\begin{remark}
显然有
\begin{align*}
D^+f(x_0)\geqslant D_+f(x_0),\quad D^-f(x_0)\geqslant D_-f(x_0), \\
D^+(-f) = -D_+(f),\quad D^-(-f) = -D_-(f).
\end{align*}
\end{remark}

\begin{example}
设 $a<b$, $a'<b'$, 作函数
\begin{align*}
f(x) =
\begin{cases}
ax\sin^2\frac{1}{x}+bx\cos^2\frac{1}{x}, & x>0, \\
0, & x=0, \\
a'x\sin^2\frac{1}{x}+b'x\cos^2\frac{1}{x}, & x<0.
\end{cases}
\end{align*}
由于在范围
\begin{align*}
\frac{1}{(2n+2)\pi}<x\leqslant\frac{1}{2n\pi}
\end{align*}
内, $\cos(1/x)$ 及 $\sin(1/x)$ 可取到 $-1$ 到 $1$ 之间的一切值, 故
\begin{align*}
D^+f(0) = \sup_{\theta}(a\sin^2\theta+b\cos^2\theta) = b.
\end{align*}
类似地, 有 $D_+f(0)=a$, $D_-f(0)=a'$, $D^-f(0)=b'$.
\end{example}

\begin{example}
设 $f\in C([a,b])$, 则存在 $x_0\in(a,b)$ 以及常数 $k$, 使得
\begin{align*}
D_-f(x_0)\geqslant k\geqslant D^+f(x_0) \quad \text{或} \quad D^-f(x_0)\leqslant k\leqslant D_+f(x_0).
\end{align*}
\end{example}
\begin{remark}
\begin{enumerate}[(1)]
\item 若 $f(x)$ 是 $[a,b]$ 上的递增函数, 则其 Dini 导(函)数可测;
\item 设
\begin{align*}
f(x) =
\begin{cases}
x, & x\in\mathbb{Q}, \\
-x, & x\notin \mathbb{Q},
\end{cases}
\quad g(x) = -f(x) \quad (x\in\mathbb{R}),
\end{align*}
则 $D^+f(0)=1$, $D^+g(0)=1$, $D^+(f+g)(0)=0$, 从而可知
\begin{align*}
D^+(f+g)(0)\neq D^+f(0)+D^+g(0).
\end{align*}
\end{enumerate}
\end{remark}
\begin{proof}
记 $k=(f(b)-f(a))/(b-a)$, 并考查 $F(x)=f(x)-kx$. 易知 $F\in C([a,b])$, 且有
\begin{align*}
F(a) &= f(a)-\frac{f(b)-f(a)}{b-a}a = \frac{1}{b-a}(bf(a)-af(b)), \\
F(b) &= \frac{1}{b-a}(bf(a)-af(b)) = F(a).
\end{align*}
由此知存在 $x_0\in(a,b)$, 使 $F(x_0)$ 是 $[a,b]$ 上 $F(x)$ 的最大值或最小值, 从而得到
\begin{align*}
D_-F(x_0)\geqslant 0,\quad D^+F(x_0)\leqslant 0 \quad (x_0 \text{ 为最大值点}), \\
D_+F(x_0)\geqslant 0,\quad D^-F(x_0)\leqslant 0 \quad (x_0 \text{ 为最小值点}).
\end{align*}
(注意, 若 $\psi'(x_0)$ 存在, 且 $g(x)=\varphi(x)+\psi(x)$, 则 $D^\pm g(x_0)=D^\pm\varphi(x_0)+\psi'(x_0)$, $D_\pm g(x)=D_\pm\varphi(x_0)+\psi'(x_0)$), 即
\begin{align*}
D_-f(x_0)\geqslant k\geqslant D^+f(x_0) \quad (x_0 \text{ 为最大值点}); \\
D_+f(x_0)\geqslant k\geqslant D^-f(x_0) \quad (x_0 \text{ 为最小值点}).
\end{align*}
\end{proof}

\begin{theorem}[Lebesgue定理(单调函数可微性定理)]\label{theorem:Lebesgue定理--实变函数}
设 $f(x)$ 为定义在 $[a,b]$ 上的单调的实值函数, 则 $f(x)$ 在 $[a,b]$ 上几乎处处可导 (可微), 且其导函数在 $[a,b]$ 上 L-可积, 并满足
\begin{align}
\int_a^b f'(x)\mathrm{d}\mu \leqslant f(b)-f(a).\label{eq:::-hr82h39rhswufh8jiojshiojfw8290rui9ejwf90382jfwds3}
\end{align}
\end{theorem}
\begin{proof}
不妨设$f$为单调递增实值函数,否则用$-f$代替即可.先证明在 $[a,b]$ 上几乎处处成立
\begin{align}
D^+f(x)=D_+f(x)=D^-f(x)=D_-f(x).\label{eq:::-hr82h39rhswufh8jiojshiojfw8290rui9ejwf90382jfwds4}
\end{align}
令 $E_1 = \{x\in(a,b): D^+f(x) > D_-f(x)\}$, 则
\begin{align*}
E_1 = \bigcup_{r,s\in\mathbb{Q}} \{x\in(a,b): D^+f(x) > r > s > D_-f(x)\}.
\end{align*}
下面证明 $m^*(E_1)=0$, 为此只需证明: 对 $\forall r,s\in\mathbb{Q}$, 都有
\begin{align*}
m^*(\{x\in(a,b): D^+f(x) > r > s > D_-f(x)\}) = 0.
\end{align*}
记
\begin{align*}
A = \{x\in(a,b): D^+f(x) > r > s > D_-f(x)\}.
\end{align*}
对任意的 $\varepsilon>0$, 存在开集 $G\supset A$, 使得 $m(G) < m^*(A)+\varepsilon$. 对任意的 $x\in A$, 由于 $D^-f(x) < s$, 故存在 $h>0$ 使得 $[x-h,x]\subset G$, 并且
\begin{align}
f(x)-f(x-h) < sh.\label{eq:::-hr82h39rhswufh8jiojshiojfw8290rui9ejwf90382jfwds5}
\end{align}
所有这样的区间 $[x-h,x]$ 构成了 $A$ 的一个Vitali覆盖. 由\hyperref[theorem:Vitali覆盖定理--实变函数]{Vitali覆盖定理}, 存在有限个互不相交的区间 $I_i = [x_i-h_i,x_i]$, $i=1,\cdots,n$, 使得 $m^*(A-\bigcup_{i=1}^n I_i) < \varepsilon$. 令 $B = A\cap\bigcup_{i=1}^n I_i^\circ$, 则
\begin{align}
m^*(A) \leqslant m^*\!\left(A\cap\bigcup_{i=1}^n I_i^\circ\right) + m^*\!\left(A-\bigcup_{i=1}^n I_i^\circ\right) < m^*(B) + \varepsilon.\label{eq:::-hr82h39rhswufh8jiojshiojfw8290rui9ejwf90382jfwds6}
\end{align}
由式 \eqref{eq:::-hr82h39rhswufh8jiojshiojfw8290rui9ejwf90382jfwds5}, 我们有
\begin{align}
\sum_{i=1}^n (f(x_i)-f(x_i-h_i)) < s\sum_{i=1}^n h_i < sm(G) < s(m^*(A)+\varepsilon).\label{eq:::-hr82h39rhswufh8jiojshiojfw8290rui9ejwf90382jfwds7}
\end{align}
对每个 $y\in B$, 由于 $D^+f(y) > r$, 故存在 $k>0$, 使得区间 $[y,y+k]$ 包含在某个区间 $I_i^\circ$ 内, 并且
\begin{align}
f(y+k)-f(y) > rk.\label{eq:::-hr82h39rhswufh8jiojshiojfw8290rui9ejwf90382jfwds8}
\end{align}
所有这样的区间 $[y,y+k]$ 构成了 $B$ 的一个Vitali覆盖. 再由\hyperref[theorem:Vitali覆盖定理--实变函数]{Vitali覆盖定理}, 存在有限个互不相交的区间 $J_j = [y_j,y_j+k_j]$, $j=1,\cdots,m$, 使得 $m^*(B-\bigcup_{j=1}^m J_j) < \varepsilon$. 利用式 \eqref{eq:::-hr82h39rhswufh8jiojshiojfw8290rui9ejwf90382jfwds6} 得
\begin{align*}
m^*(A)&<m^*(B)+\varepsilon \leqslant m^*\!\left( B\cap \bigcup_{j=1}^m{J_{j}^{\circ}} \right) +m^*\!\left( B-\bigcup_{j=1}^m{J_{j}^{\circ}} \right) +\varepsilon 
\\
&\leqslant m\!\left( \bigcup_{j=1}^m{J_j} \right) +2\varepsilon =\sum_{j=1}^m{k_j}+2\varepsilon .
\end{align*}
因此, $\sum_{j=1}^m k_j > m^*(A)-2\varepsilon$. 再由式 \eqref{eq:::-hr82h39rhswufh8jiojshiojfw8290rui9ejwf90382jfwds8}, 我们有
\begin{align}
\sum_{j=1}^m (f(y_j+k_j)-f(y_j)) > r\sum_{j=1}^m k_j > r(m^*(A)-2\varepsilon).\label{eq:::-hr82h39rhswufh8jiojshiojfw8290rui9ejwf90382jfwds9}
\end{align}
由于 $f(x)$ 是单调递增函数, 并且每个 $J_j$ 包含在某个 $I_i$ 中, 于是
\begin{align}
\sum_{i=1}^n (f(x_i)-f(x_i-h_i)) \geqslant \sum_{j=1}^m (f(y_j+k_j)-f(y_j)).\label{eq:::-hr82h39rhswufh8jiojshiojfw8290rui9ejwf90382jfwds10}
\end{align}
结合式 \eqref{eq:::-hr82h39rhswufh8jiojshiojfw8290rui9ejwf90382jfwds7}, \eqref{eq:::-hr82h39rhswufh8jiojshiojfw8290rui9ejwf90382jfwds9} 和 \eqref{eq:::-hr82h39rhswufh8jiojshiojfw8290rui9ejwf90382jfwds10}, 可得
\begin{align*}
r(m^*(A)-2\varepsilon) < s(m^*(A)+\varepsilon).
\end{align*}
由 $\varepsilon$ 的任意性得 $rm^*(A)\leqslant sm^*(A)$. 由于 $r>s$, 故必有 $m^*(A)=0$. 由此得到 $m^*(E_1)=0$. 类似地, 若令
\begin{align*}
E_2 = \{x\in(a,b): D^-f(x) > D_+f(x)\},
\end{align*}
则可以证明 $m^*(E_2)=0$, 令 $E=E_1\cup E_2$, 则 $m^*(E)=0$. 在 $(a,b)-E$ 上, 我们有
\begin{align*}
D_+f(x)\leqslant D^+f(x)\leqslant D_-f(x)\leqslant D^-f(x)\leqslant D_+f(x).
\end{align*}
因此, 在 $(a,b)-E$ 上式 \eqref{eq:::-hr82h39rhswufh8jiojshiojfw8290rui9ejwf90382jfwds4} 成立. 这表明极限
\begin{align*}
g(x) = \lim_{h\to 0} \frac{f(x+h)-f(x)}{h}
\end{align*}
几乎处处存在 (有限或 $\pm\infty$). 当 $g(x)$ 有限时, $f$ 在点 $x$ 可导. 令
\begin{align*}
g_n(x) = n\left[f\!\left(x+\frac{1}{n}\right)-f(x)\right],\quad n=1,2,\cdots,
\end{align*}
其中定义当 $x>b$ 时, $f(x)=f(b)$, 则 $g_n\xrightarrow{\text{a.e.}} g$, 从而 $g$ 可测. 又 $f$ 单调递增, 故 $g_n$ 非负. 于是
\begin{align*}
\int_a^b g_n(x)\mathrm{d}\mu &= n\int_a^b \left[f\!\left(x+\frac{1}{n}\right)-f(x)\right]\mathrm{d}\mu \\
&= n\int_{a+\frac{1}{n}}^{b+\frac{1}{n}} f(x)\mathrm{d}\mu - n\int_a^b f(x)\mathrm{d}\mu \\
&= n\int_b^{b+\frac{1}{n}} f(x)\mathrm{d}\mu - n\int_a^{a+\frac{1}{n}} f(x)\mathrm{d}\mu \\
&= f(b)-n\int_a^{a+\frac{1}{n}} f(x)\mathrm{d}\mu.
\end{align*}
因此, 由\hyperref[theorem:Fatou引理]{Fatou引理}可得
\begin{align}
\int_a^b g(x)\mathrm{d}\mu \leqslant \liminf_{n\to\infty} \int_a^b g_n(x)\mathrm{d}\mu  = \liminf_{n\to\infty} \left(f(b)-n\int_a^{a+\frac{1}{n}} f(x)\mathrm{d}\mu\right) \leqslant f(b)-f(a). \label{eq:::iojfw8}
\end{align}
这表明 $g(x)$ 是 L-可积的, 因此 $g(x)$ 几乎处处有限. 于是, $f$ 几乎处处可导. 由于 $f'(x)=g(x)$, a.e. $x\in[a,b]$, 故式 \eqref{eq:::iojfw8} 表明式 \eqref{eq:::-hr82h39rhswufh8jiojshiojfw8290rui9ejwf90382jfwds3} 成立.
\end{proof}

\begin{example}
设 $E$ 为 $(a,b)$ 中的零测集, 对每个 $n\in\mathbb{N}$, 取开集 $G_n\supset E$ 满足 $m(G_n)<\frac{1}{2}^n$. 令
\begin{align*}
f_n(x) = m((a,x)\cap G_n),\quad x\in(a,b),\quad n=1,2,\cdots
\end{align*}
显然, $f_n(x)$ 单调递增, 非负, 且满足 $|f_n|\leqslant \frac{1}{2}^n$. 又对 $\forall x\in(a,b)$, 当 $a+h\in(a,b)$ 时
\begin{align*}
f(x+h)-f(x)\leqslant |h|,
\end{align*}
故 $f_n(x)$ 是 $(a,b)$ 上的连续函数.

令 $f(x)=\sum_{n=1}^\infty f_n(x)$, 则 $f(x)$ 为单调递增的非负连续函数. 现设 $x\in E$, 对 $\forall n\in\mathbb{N}$, 当 $|h|$ 充分小时, 有 $[x,x+h]\subset G_i\cap(a,b)$, $i=1,\cdots,n$. 由 $f_n$ 的定义, 此时有
\begin{align*}
\frac{f_i(x+h)-f_i(x)}{h}=1,\quad i=1,\cdots,n
\end{align*}
于是
\begin{align*}
\frac{f(x+h)-f(x)}{h}\geqslant\sum_{i=1}^n\frac{f_i(x+h)-f_i(x)}{h}=n.
\end{align*}
因此, $f'(x)=+\infty$. 这表明 $f(x)$ 在 $E$ 上处处不可导.
\end{example}
\begin{remark}
这个例题表明,\hyperref[theorem:Lebesgue定理--实变函数]{Lebesgue定理}中单调函数是几乎处处可微的这一结论,一般来说是不能改进的.
\end{remark}

\begin{theorem}[Fubini 逐项微分定理]\label{theorem:实变函数--Fubini 逐项微分定理}
设 $\{f_n(x)\}$ 是 $[a,b]$ 上的递增函数列, 且 $\sum_{n=1}^\infty f_n(x)$ 在 $[a,b]$ 上收敛, 则
\begin{align*}
\frac{\mathrm{d}}{\mathrm{d}x}\left(\sum_{n=1}^\infty f_n(x)\right) = \sum_{n=1}^\infty \frac{\mathrm{d}}{\mathrm{d}x}f_n(x), \quad \text{a.e. } x\in[a,b].
\end{align*}
\end{theorem}
\begin{remark}
联系到微积分中所学的逐项微分定理, 上例部分地反映 Lebesgue 积分理论的优越性, 只不过这里只有几乎处处成立的结论. 这一点是无法克服的, 即使每个 $f_n(x)$ 都是可微的递增函数, 其和函数也可微, 也不保证对一切 $x$, 有
\begin{align*}
\frac{\mathrm{d}}{\mathrm{d}x}\left(\sum_{n=1}^\infty f_n(x)\right) = \sum_{n=1}^\infty \frac{\mathrm{d}}{\mathrm{d}x}f_n(x).
\end{align*}
\end{remark}
\begin{proof}
首先, $\sum_{n=1}^\infty f_n(x)$ 是 $[a,b]$ 上的递增函数, 故它是几乎处处可微的. 其次, 由于
\begin{align*}
\sum_{n=1}^\infty f_n(x) = \sum_{n=1}^N f_n(x) + R_N(x), \quad R_N(x) = \sum_{n=N+1}^\infty f_n(x)
\end{align*}
($R_N(x)$ 也是几乎处处可微的), 我们有
\begin{align*}
\frac{\mathrm{d}}{\mathrm{d}x}\left(\sum_{n=1}^\infty f_n(x)\right) = \sum_{n=1}^N \frac{\mathrm{d}}{\mathrm{d}x}f_n(x) + \frac{\mathrm{d}}{\mathrm{d}x}R_N(x), \quad \text{a.e. } x\in[a,b].
\end{align*}
从而只需指出
\begin{align*}
\lim_{N\to\infty} R_N'(x) = 0, \quad \text{a.e. } x\in[a,b]
\end{align*}
即可. 注意到 $f_n'(x)\geqslant 0$ ($n=1,2,\cdots$, a.e. $x\in[a,b]$), 故有
\begin{align*}
R_N'(x) = f_{N+1}'(x) + R_{N+1}'(x) \geqslant R_{N+1}'(x), \quad \text{a.e. } x\in[a,b].
\end{align*}
这说明存在
\begin{align*}
\lim_{N\to\infty} R_N'(x) \stackrel{\text{def}}{=} R(x) \geqslant 0, \quad \text{a.e. } x\in[a,b].
\end{align*}
根据公式\hyperref[theorem:Lebesgue定理--实变函数]{Lebesgue定理}, 可知
\begin{align*}
\int_a^b \lim_{N\to\infty} R_N'(x)\mathrm{d}x = \lim_{N\to\infty} \int_a^b R_N'(x)\mathrm{d}x 
\leqslant \lim_{N\to\infty} (R_N(b)-R_N(a)) = 0.
\end{align*}
由此即得 $\lim_{N\to\infty} R_N'(x) = 0$, a.e. $x\in[a,b]$.
\end{proof}

\begin{example}
记 $[0,1]$ 中的有理数全体为 $\{r_n\}$, 定义
\begin{align*}
f_n(x) =
\begin{cases}
0, & 0\leqslant x < r_n, \\
\frac{1}{2^n}, & r_n\leqslant x\leqslant 1,
\end{cases}
\end{align*}
其中 $n=1,2,\cdots$, 则 $f(x)$ 在 $[0,1]$ 上单调递增, 显然 $f_n'(x)=0$, a.e. $x\in[0,1]$. 令
\begin{align*}
f(x) = \sum_{n=1}^\infty f_n(x),\quad x\in[0,1].
\end{align*}
注意到
\begin{align*}
f(x) = \sum_{r_n\leqslant x} \frac{1}{2^n},\quad 0<x<1.
\end{align*}
易知 $f(x)$ 在 $[0,1]$ 上严格递增, 由\hyperref[theorem:实变函数--Fubini 逐项微分定理]{Fubini 逐项微分定理}得, $f'(x)=\sum_{n=1}^\infty f_n'(x)=0$, a.e. $x\in[0,1]$. 于是
\begin{align*}
\int_0^1 f'(x)\mathrm{d}\mu = 0 < f(1)-f(0)
\end{align*}
(其中 $f(1)=1$). 这说明, 一般情形下\hyperref[theorem:Lebesgue定理--实变函数]{Lebesgue定理}中的不等式\eqref{eq:::-hr82h39rhswufh8jiojshiojfw8290rui9ejwf90382jfwds3}不能成为等式.
\end{example}
\begin{remark}
这个例题就是利用\hyperref[theorem:实变函数--Fubini 逐项微分定理]{Fubini 逐项微分定理}构造一个实例来说明一般情形下,\hyperref[theorem:Lebesgue定理--实变函数]{Lebesgue定理}中的不等式\eqref{eq:::-hr82h39rhswufh8jiojshiojfw8290rui9ejwf90382jfwds3}不能成为等式.
\end{remark}

\begin{example}
设 $f\in C((-\infty,+\infty))$, 且令
\begin{align*}
f_t(x) = f(x+t)-f(x), \quad -\infty < x, t < +\infty.
\end{align*}
若对任意的 $t\in(-\infty,+\infty)$, $f_t(x)$ 对 $x\in(-\infty,+\infty)$ 可微, 则 $f(x)$ 在 $(-\infty,+\infty)$ 上可微.
\end{example}
\begin{proof}
\begin{enumerate}[(1)]
\item 对任意的两点 $x'$ 和 $x''=x'+t$ ($-\infty<t<+\infty$), 我们有等式
\begin{align*}
\frac{f(x''+h)-f(x'')}{h} &= \frac{f(x'+t+h)-f(x'+t)}{h} \\
&= \frac{f(x'+t+h)-f(x'+h)-[f(x'+t)-f(x')]}{h} \\
&\quad + \frac{f(x'+h)-f(x')}{h} \\
&= \frac{f_t(x'+h)-f_t(x')}{h} + \frac{f(x'+h)-f(x')}{h}.
\end{align*}
由此和题设可知
\begin{align*}
D^\pm_\pm f(x'') = f_t'(x') + D^\pm_\pm f(x'),
\end{align*}
且只要 $f(x)$ 在一点可微, 就可知 $f(x)$ 处处可微了.

\item 根据例 3, 可知存在 $k$ 以及 $x_1$, 使得(不妨假定)
\begin{align*}
D_-f(x_1) \geqslant k \geqslant D^+f(x_1).
\end{align*}
从而对任意的 $x=x_1+t$ ($-\infty<t<+\infty$), 有
\begin{align*}
D_-f(x) = f_t'(x_1) + D_-f(x_1) \geqslant f_t'(x_1) + D^+f(x_1) = D^+f(x).
\end{align*}

\item 不妨假定 $f(x)$ 不是上凸函数(否则就有可微点了), 则存在区间 $[a,b]$, 使得在 $[a,b]$ 上 $f(x)$ 位于点 $(a,f(a))$ 与点 $(b,f(b))$ 连接线的下方. 现在记
\begin{align*}
F(x) = f(x)-lx, \quad l = \frac{f(b)-f(a)}{b-a},
\end{align*}
则易知存在 $x_2\in(a,b)$, 使得 $F(x)$ 在 $x=x_2$ 处取得最小值. 由此得
\begin{align*}
D^-F(x_2) \leqslant 0, \quad &\text{即} \quad D^-f(x_2) \leqslant l; \\
D_+F(x_2) \geqslant 0, \quad &\text{即} \quad D_+f(x_2) \geqslant l.
\end{align*}
从而对任意的 $x\in(-\infty,+\infty)$, 又有
\begin{align*}
D^-f(x) \leqslant D_+f(x).
\end{align*}

\item 综合上述结论, 我们有
\begin{align*}
D_-f(x) = D^-f(x) = D_+f(x) = D^+f(x), \quad -\infty < x < +\infty.
\end{align*}
这说明 $f'(x)$ 有意义, 从而 $f(x)$ 就有可微点了.

\end{enumerate}
\end{proof}




































\end{document}